\section{Evaluation Setting and Outcome-Level Landscape}
\label{sec:main_exp}

We first evaluate seven frontier models on the same tasks with a shared harness and protocol to establish a controlled comparison of final performance and resource use.
This outcome-level landscape anchors our subsequent analyses of research behavior, experience reuse, harness effects, and solution novelty.

\subsection{Evaluation Setting}
\label{sec:setup}

\textbf{Evaluation tasks.}
Our evaluation focuses on four workload families that capture distinct demands of automated research: \textit{Model Development}, \textit{System Optimization}, \textit{Puzzle \& Challenge}, and \textit{CUDA}.
We instantiate this scope with 36 expert-curated tasks from AutoLab~\citep{xu2026autolab}, comprising 7, 15, 10, and 4 tasks from the four families, respectively.
Each task provides an objective, a correct but deliberately suboptimal starting artifact, an expert-written reference solution, a wall-clock budget, and an automated verifier.
Within the allotted budget, the agent iteratively improves the artifact, and the verifier scores the final submission relative to the starting artifact and expert reference on a normalized scale from $0$ to $1$.

\textbf{Models and harness.}
We evaluate seven frontier models: \textbf{Claude-Opus-4.7}~\citep{anthropic2026opus47}, \textbf{GPT-5.5}~\citep{openai2026gpt55}, \textbf{Gemini-3.1-Pro}~\citep{google2026gemini31pro}, \textbf{GLM-5.2}~\citep{zai2026glm5.2}, \textbf{Kimi-K2.7-Code}~\citep{moonshot2026kimik27code}, \textbf{DeepSeek-V4-Pro}~\citep{deepseek2026deepseekv4}, and \textbf{LongCat-2.0}~\citep{longcat2026longcat2}.\footnote{These were generally the latest available model versions from their respective providers when this study was initiated. All evaluations were conducted from June~2026 to July~10,~2026, using the provider API versions available at the time.}
For the main cross-model comparison, all models use \textbf{Claude Code} (v2.1.152) as a practical shared harness, thereby holding the tool interface and iteration policy fixed.
We separately evaluate how harness choice affects performance by comparing Claude Code with model-native and open-source alternatives in \S\ref{sec:harness_ablation}.


\textbf{Evaluation protocol and metrics.}
Given the open-ended, long-horizon nature of auto research and the resulting variation across runs, we evaluate each model on all 36 tasks with three independent rollouts per model--task pair, yielding 756 rollouts.
For each three-rollout set, we report \textbf{avg@3} and \textbf{best@3} to characterize the model's typical and best-observed performance, respectively.
Each task retains its original wall-clock budget of 2--12 hours, determined by workload scale.
To enable later process analysis, we add only a record-keeping instruction requiring the agent to commit after each iteration and maintain an experiment journal; all other task and execution conditions remain unchanged.
A complete task instruction is provided in Appendix~\ref{app:task_instruction_example}.


\subsection{Outcome-Level Results}
\label{sec:main_results}

\begin{figure}
    \centering
    \includegraphics[width=\linewidth]{Figures/main_outcome_landscape.pdf}
    \caption{Outcome-level performance across seven models.
    Solid segments indicate avg@3, while full bar heights indicate best@3. (a) Overall performance, ranked by avg@3.
    (b) Category-level performance; filled and open circles mark the avg@3 and best@3 leaders, respectively.}
    \label{fig:main_exp_results}
\end{figure}

\textbf{Overall performance.}
Figure~\ref{fig:main_exp_results}(a) reveals a clear overall hierarchy.
Opus-4.7 ranks first on both avg@3 ($0.739$) and best@3 ($0.790$), combining the strongest average performance with the highest observed ceiling.
GPT-5.5, GLM-5.2, and Gemini-3.1-Pro form a compact second tier, spanning only $0.029$ on avg@3 and $0.022$ on best@3.
Within this tier, GPT exhibits the highest performance ceiling, whereas GLM delivers the strongest stable performance across runs.

\textbf{Average performance separates models more sharply than best performance.}
Across models, the highest-to-lowest gap is $0.237$ under avg@3 but only $0.122$ under best@3.
For example, Kimi's best@3 is only $0.028$ below GLM and $0.021$ below Gemini, but it falls substantially farther behind both models on avg@3.
Lower-ranked models can therefore reach competitive solutions, but do so less consistently across repeated runs.
Our subsequent analyses show that harness design and experience reuse can help narrow this consistency gap, while training objectives based on relative outcomes across repeated rollouts offer a complementary direction.

\textbf{Task categories reveal distinct capability profiles.}
Figure~\ref{fig:main_exp_results}(b) shows that Opus's overall lead is broad but not universal: it leads avg@3 on Model Development, System Optimization, and CUDA, whereas GLM narrowly leads Puzzle \& Challenge.
Puzzle \& Challenge appears the most accessible category, producing high scores across models and the smallest highest-to-lowest gaps ($0.150$ on avg@3 and $0.074$ on best@3).
By contrast, CUDA is both lower-scoring and the most separating, with corresponding gaps of $0.403$ and $0.414$, indicating that low-level GPU optimization remains substantially more difficult.
CUDA also reveals different strengths under the two metrics: Opus leads avg@3, whereas GPT leads best@3, indicating that GPT can reach stronger solutions but does so less consistently.
Category-level evaluation therefore reveals workload-specific strengths and differences in reliability that an overall score cannot capture.
The full category-level breakdown is reported in Appendix~\ref{app:outcome_by_category}.


\begin{figure}[t]
    \centering
    \includegraphics[width=\linewidth]{Figures/task_type_mean_cost.pdf}
    \caption{Mean estimated inference cost per task across four task categories and overall. Values average over the three independent rollouts for each model--task pair. For consistent cross-model comparison, all input tokens are priced without cache discounts.}
    \label{fig:cost_by_task_type}
\end{figure}

\subsection{Cost and Resource Analysis}
We record token consumption and wall-clock time for the main evaluation runs, and use public API prices to estimate each model's mean inference cost per task.
Figure~\ref{fig:cost_by_task_type} reports mean cost per task for each workload family and overall.
Together with the performance results in Figure~\ref{fig:main_exp_results}, the cost view shows that Opus-4.7 achieves the strongest best@3 ($0.790$), but at a much higher mean cost of \$89.9 per task; GPT-5.5 and GLM-5.2 provide close alternatives ($0.772$ and $0.757$) for substantially less (\$16.5 and \$33.0 per task).
LongCat-2.0 and DeepSeek-V4-Pro trade some performance for very low mean costs (\$3.9 and \$4.3 per task), making them attractive options under tight budgets.
By category, CUDA tasks are the most expensive on average, while Puzzle \& Challenge tasks are the cheapest, a pattern broadly consistent with their relative difficulty.
Appendix~\ref{app:resource_by_category} further reports the wall-clock time and token consumption, and analyzes the relationship between overall performance and resource use.

